\vspace{-3mm}
\section{Conclusion and Limitations} \label{sec:conclusion}
\vspace{-2mm}
We presented FisherRF, a novel method for active view selection, active mapping, and uncertainty quantification in Radiance Fields. 
Leveraging Fisher Information, our method provides an efficient and effective means to quantify the observed information of Radiance Field models. The flexibility of our approach allows it to be applied to various model parametrizations, including 3D Gaussian Splatting and Plenoxels.
Our extensive evaluation of active view selection, active mapping, and uncertainty quantification has consistently shown superior performance compared to existing methods and heuristic baselines. These results highlight the potential of our approach to significantly enhance the quality and efficiency of image rendering and reconstruction tasks with limited viewpoints.
However, our method is limited to static scenes in a confined scenario. Reconstructing dynamically changing the Radiance Field and quantifying its Fisher Information is still an open problem.
More work could be done to overcome the limitations and extend the proposed method to more challenging settings. 

{\small
\paragraph{Acknowledgements}
The authors gratefully appreciate support through the following grants: NSF FRR 2220868, NSF IIS-RI 2212433, NSF TRIPODS 1934960, NSF CPS 2038873.
The authors thank Pratik Chaudhari for the insightful discussion and Yinshuang Xu for proofreading the drafts.
}

